\subsection{Fragment Invariance and Reuse}
\label{sec:invariance}

The compiled fragment is determined by $Z_{\mathrm{ref}}$. A witnessing
homomorphism also depends on the implementation and on both encodings. This
difference determines what can be reused when a design changes.

Representational changes over a system lifecycle should not invalidate an
established behavioral result.

\begin{theorem}[Encoding invariance]
\label{thm:encoding-invariance}
Replacing either $Z_{\mathrm{ref}}$ or $Z_{\mathrm{impl}}$ by an isomorphic
system preserves satisfaction of the compiled dynamics fragment. The same
holds for T3SD copies that rename or permute ports.
\end{theorem}

\begin{proof}[Proof sketch]
Compose the witnessing homomorphism with the isomorphism in the required
direction. The inverse isomorphism proves the reverse implication. T3SD copies
carry the corresponding bijections on state and boundary indices, so the same
composition argument applies.
\end{proof}

Component-level evidence must lift to assembled resultants if verification is
to scale with system architecture.

\begin{theorem}[Elaboration and coupling]
\label{thm:elaboration-coupling}
Dynamics-fragment conformance composes through successive homomorphic
realisations on either side. For two coupling recipes with the same free
boundary, port-preserving component homomorphisms induce a homomorphism between
their resultants.
\end{theorem}

\begin{proof}[Proof sketch]
Composition of the state, input, and output maps preserves surjectivity and the
homomorphism equations. For coupled systems, take the componentwise product of
the state maps. Port preservation makes the internal-wire values agree, so the
product map preserves the resultant transition and readout functions. Standard
T3SD coupling results supply the corresponding resultant construction
\citep{Wymore1993}.
\end{proof}

Reindexing components, flattening nested recipes, and removing
boundary-irrelevant components are structural instances of
Theorem~\ref{thm:elaboration-coupling}. Section~\ref{sec:case-machine-lift}
illustrates component elaboration by changing the instruction table's phase
encoding.

Together, these results replace one monolithic obligation with a hierarchy of checks.
Each component can be checked against its reference fragment; the coupling
theorem lifts the component results when the required coupling witness is
supplied; subsequent elaborations reuse the established fragment. Finite
components admit propositional checking, while infinite components may use
symbolic proofs or finite abstractions. The RLE case proves conformance of the
full resultant directly and illustrates fragment reuse on its instruction
table. The underlying search problem remains difficult in general, while
decomposition and reuse provide a route to larger architectures.

A characteristic fragment should distinguish finite behavior, not merely admit
coarse mutual projections.

\begin{theorem}[Finite canonicity]
\label{thm:finite-canonicity}
Let $Z_1$ and $Z_2$ be finite systems. If each satisfies the dynamics fragment
compiled from the other, then $Z_1$ and $Z_2$ are isomorphic.
\end{theorem}

\begin{proof}[Proof sketch]
Mutual satisfaction yields surjective homomorphisms in both directions by
Theorem~\ref{thm:dynamics-characterization}. Finiteness forces the state,
input, and output maps to be bijective. Their inverses preserve the same
transition and readout equations, giving an isomorphism.
\end{proof}

Theorems~\ref{thm:encoding-invariance}--\ref{thm:finite-canonicity} support
incremental verification while preserving Wymore's implementation semantics.
Composition lifts established component results, and finite presentations turn
local conformance obligations into generated decision problems.
